\subsection{局部可忽略的集与函数}

定义 6. --- 称函数 \(f \in \mathcal{F}_{+}\) （相应地 T 的子集 A）对测度 \(\mu\) 是局部可忽略的，如果 \(\mu^{\bullet}(f) = 0\) （相应地 \(\mu^{\bullet}(A) = 0\) ）。称 \(\mu\) 集中于 T 的子集 A 上，如果 T - A 局部 \(\mu\) -可忽略。

注记. --- 1) 当 T 局部紧时，如此定义的概念与通常的概念一致。

\begin{enumerate}
\def\labelenumi{\arabic{enumi})}
\setcounter{enumi}{1}
\item
  在定义过可忽略集之后，我们将看到，局部可忽略集正是在 T 的每一点处的芽都是某个可忽略集之芽的那些集合 (No.~9, Prop. 14 的 Cor. 2)。
\item
  如同 Chs. IV 与 V 中那样，“局部几乎处处”一语与“除去一个局部可忽略集外”同义。
\item
  若 \(\theta\) 是一个复测度，我们说一个函数（相应地 T 的一个子集）对 \(\theta\) 局部可忽略，如果它对正测度 \(|\theta|\) 局部可忽略。
\end{enumerate}

例. --- 设 L 是 T 的紧子集，\(\lambda\) 是 L 上的一个测度，\(\mu\) 是由 \(\lambda\) 定义的 T 上的测度 (No.~3, 例 2)。公式 (3) 立即表明：函数 \(f \in \mathcal{F}_{+}(\mathrm{T})\) 局部 \(\mu\) -可忽略当且仅当 \(f_{L}\) 是 \(\lambda\) -可忽略的。

由公式 (1) 立即得出：函数 \(f \in \mathcal{F}_{+}(\mathrm{T})\) 局部 \(\mu\) -可忽略当且仅当对 T 的每个紧子集 K，\(f_{K}\) 是 \(\mu_{K}\) -可忽略的。于是局部可忽略集的性质立即化归为紧空间中可忽略集的性质，后者已在 Ch. IV 中处理。下面是一些今后将直接引用的结果。

--- 为使函数 \(f \geqslant 0\) 局部可忽略，必须且只需 \(f(t) = 0\) 局部几乎处处成立 (Ch. IV, §2, No.~3, Th. 1)。若 f 是以 Banach 空间为值的函数，则说 f = 0 局部几乎处处成立与说 \(\mu^{\bullet}(|\mathbf{f}|) = 0\) 等价；这时我们仍称 f 是局部可忽略的。

--- 局部可忽略函数 \(\geqslant0\) 的序列的和与上包络都是局部可忽略的 (loc. cit., No.~1, Prop. 2)。

--- 若 \(f\) 与 \(g\) 是两个 \(\geqslant 0\) 的函数且局部几乎处处相等，则 \(\mu^{\bullet}(f) = \mu^{\bullet}(g)\) (loc. cit., No.~3, Prop. 6)。

